% sec4_1.tex — Section 4.1: The Multiple-Equation Model

\section{The Multiple-Equation Model}
\label{sec:4.1}

To be able to deal with more than one equation, we need to complicate the notation
by introducing an additional subscript designating the equation. For example, the
dependent variable of the $m$-th equation for observation $i$ will be denoted $y_{im}$.

\subsection*{Linearity}

There are $M$ linear equations, each of which is a linear equation like the one in
Assumption~3.1:

\begin{assumption}[linearity]
\label{ass:4.1}
\textit{There are $M$ linear equations}
\begin{equation}
y_{im} = \mathbf{z}_{im}' \boldsymbol{\delta}_m + \varepsilon_{im}
\quad (m = 1, 2, \ldots, M;\, i = 1, 2, \ldots, n),
\label{eq:4.1.1}
\end{equation}
\textit{where $n$ is the sample size, $\mathbf{z}_{im}$ is the $L_m$-dimensional vector of regressors,
$\boldsymbol{\delta}_m$ is the conformable coefficient vector, and $\varepsilon_{im}$ is an unobservable
error term in the $m$-th equation.}
\end{assumption}

The model will make no assumptions about the \textbf{interequation} (or \textbf{contemporaneous})
\textbf{correlation} between the errors $(\varepsilon_{i1}, \ldots, \varepsilon_{iM})$.
Also, no \textit{a priori} restrictions are placed on the coefficients from different equations.
That is, the model assumes no \textbf{cross-equation restrictions} on the coefficients.
These points can be illustrated in

\begin{quote}
\textbf{Example 4.1 (wage equation):}\quad
In the Griliches exercise of Chapter~3, we estimated the wage equation.
In one specification, Griliches adds to it the equation for $KWW$ (score on the
``Knowledge of the World of Work'' test):
\begin{align*}
LW_i &= \phi_1 + \beta_1 S_i + \gamma_1 IQ_i + \pi\, EXPR_i + \varepsilon_{i1}, \\
KWW_i &= \phi_2 + \beta_2 S_i + \gamma_2 IQ_i + \varepsilon_{i2},
\end{align*}
where (to refresh your memory) $LW_i$ is the log wage of the $i$-th individual in
the first year of the survey, $S_i$ is schooling in years, $IQ_i$ is IQ, and $EXPR_i$ is
experience in years. In this example, $M = 2$, $L_1 = 4$, $L_2 = 3$,
$y_{i1} = LW_i$, $y_{i2} = KWW_i$, and
\begin{equation}
\mathbf{z}_{i1} = \begin{bmatrix} 1 \\ S_i \\ IQ_i \\ EXPR_i \end{bmatrix}, \quad
\mathbf{z}_{i2} = \begin{bmatrix} 1 \\ S_i \\ IQ_i \end{bmatrix}, \quad
\boldsymbol{\delta}_1 = \begin{bmatrix} \phi_1 \\ \beta_1 \\ \gamma_1 \\ \pi \end{bmatrix}, \quad
\boldsymbol{\delta}_2 = \begin{bmatrix} \phi_2 \\ \beta_2 \\ \gamma_2 \end{bmatrix}.
\label{eq:4.1.2}
\end{equation}
\end{quote}

The errors $\varepsilon_{i1}$ and $\varepsilon_{i2}$ can be correlated. The correlation arises if, for example,
there is an unobservable individual characteristic that affects both the
wage rate and the test score. There are no cross-equation restrictions such as
$\beta_1 = \beta_2$.

Cross-equation restrictions naturally arise in \textbf{panel data} models where the
same relationship can be estimated for different points in time.

\begin{quote}
\textbf{Example 4.2 (wage equation for two years):}\quad
The NLS-Y data we used for the Griliches exercise actually has information for 1980,
so we can estimate two wage equations, one for the first point in time (say, 1969)
and the other for 1980:
\begin{align*}
LW69_i &= \phi_1 + \beta_1 S69_i + \gamma_1 IQ_i + \pi_1 EXPR69_i + \varepsilon_{i1}, \\
LW80_i &= \phi_2 + \beta_2 S80_i + \gamma_2 IQ_i + \pi_2 EXPR80_i + \varepsilon_{i2}.
\end{align*}
In the language of multiple equations, $S69_i$ (education in 1969) and $S80_i$
(education in 1980) are two different variables. It would be natural to entertain
the case where the coefficients remained unchanged over time, with
$\phi_1 = \phi_2$, $\beta_1 = \beta_2$, $\gamma_1 = \gamma_2$, $\pi_1 = \pi_2$.
This is a set of linear cross-equation restrictions.
\end{quote}

As will be shown in Section~4.3 below, testing (linear or nonlinear) cross-equation
restrictions is straightforward. Estimation while imposing the restriction that the
coefficients are the same across equations will be considered in Section~4.6.

\subsection*{Stationarity and Ergodicity}

Let $\mathbf{x}_{im}$ be the vector of $K_m$ instruments for the $m$-th equation. Although not
frequently the case in practice, the set of instruments can differ across equations,
and so the number of instruments, $K_m$, can depend on the equation.

\begin{quote}
\textbf{Example 4.1 (continued):}\quad
Suppose that $IQ_i$ is endogenous in both equations
and that there is a variable, $MED_i$ (mother's education), that is predetermined
in the first equation (so $\E(MED_i\, \varepsilon_{i1}) = 0$). So the instruments for the $LW$
equation are $(1, S_i, EXPR_i, MED_i)$. If these instruments are also orthogonal
to $\varepsilon_{i2}$, then the $KWW$ equation can be estimated with the same set of instruments
as the $LW$ equation. In this example, $K_1 = K_2 = 4$, and
\begin{equation}
\mathbf{x}_{i1} = \mathbf{x}_{i2} = \begin{bmatrix} 1 \\ S_i \\ EXPR_i \\ MED_i \end{bmatrix}.
\label{eq:4.1.3}
\end{equation}
\end{quote}

The multiple-equation version of Assumption~3.2 is

\begin{assumption}[ergodic stationarity]
\label{ass:4.2}
\textit{Let\/ $\mathbf{w}_i$ be unique and nonconstant elements
of $(y_{i1}, \ldots, y_{iM}, \mathbf{z}_{i1}, \ldots, \mathbf{z}_{iM},
\mathbf{x}_{i1}, \ldots, \mathbf{x}_{iM})$.
$\{\mathbf{w}_i\}$ is jointly stationary and ergodic.}
\end{assumption}

This assumption is stronger than just assuming that ergodic stationarity is satisfied
for each equation of the system. Even if $\{y_{im}, \mathbf{z}_{im}, \mathbf{x}_{im}\}$
is stationary and ergodic for each individual equation $m$, it does not necessarily follow,
strictly speaking, that the larger process $\{\mathbf{w}_i\}$, which is the union of individual
processes, is (jointly) stationary and ergodic (see Example~2.3 of Chapter~2 for an
example of a vector nonstationary process whose components are univariate stationary
processes). In practice, the distinction is somewhat blurred because the equations
often share common variables.

\begin{quote}
\textbf{Example 4.1 (continued):}\quad
Because $\mathbf{z}_{i1}$, $\mathbf{z}_{i2}$, $\mathbf{x}_{i1}$,
and $\mathbf{x}_{i2}$ for the example have some common elements,
$\mathbf{w}_i$ has only six elements:
$(LW_i, KWW_i, S_i, IQ_i, EXPR_i, MED_i)$.
The ergodic stationarity assumption for the first equation
of this example is that $\{LW_i, S_i, IQ_i, EXPR_i, MED_i\}$ is jointly stationary and
ergodic. Assumption~4.2 requires that the joint process include $KWW_i$.
\end{quote}

\subsection*{Orthogonality Conditions}

The orthogonality conditions for the $M$-equation system are just a collection of the
orthogonality conditions for individual equations.

\begin{assumption}[orthogonality conditions]
\label{ass:4.3}
\textit{For each equation $m$, the $K_m$ variables in $\mathbf{x}_{im}$ are predetermined.
That is, the following orthogonality conditions are satisfied:}
\[
\E(\mathbf{x}_{im} \cdot \varepsilon_{im}) = \mathbf{0}
\quad (m = 1, 2, \ldots, M).
\]
\end{assumption}

So there are $\sum_m K_m$ orthogonality conditions in total. If we define
\begin{equation}
\underset{(\sum_{m=1}^{M} K_m \times 1)}{\mathbf{g}_i}
\equiv \begin{bmatrix} \mathbf{x}_{i1} \cdot \varepsilon_{i1} \\ \vdots \\
\mathbf{x}_{iM} \cdot \varepsilon_{iM} \end{bmatrix},
\label{eq:4.1.4}
\end{equation}
then all those orthogonality conditions can be written compactly as
$\E(\mathbf{g}_i) = \mathbf{0}$.
Note that we are not assuming ``cross'' orthogonalities; for example, $\mathbf{x}_{i1}$
does not have to be orthogonal to $\varepsilon_{i2}$, although it is required to be orthogonal
to $\varepsilon_{i1}$. However, if a variable is included in both $\mathbf{x}_{i1}$
and $\mathbf{x}_{i2}$, then the assumption does imply that the variable is orthogonal to
both $\varepsilon_{i1}$ and $\varepsilon_{i2}$.

\begin{quote}
\textbf{Example 4.1 (continued):}\quad
$\sum_m K_m$ in the current example is $8\; (= 4 + 4)$
and $\mathbf{g}_i$ is
\[
\mathbf{g}_i = \begin{bmatrix}
\varepsilon_{i1} \\ S_i\, \varepsilon_{i1} \\ EXPR_i\, \varepsilon_{i1} \\
MED_i\, \varepsilon_{i1} \\
\varepsilon_{i2} \\ S_i\, \varepsilon_{i2} \\ EXPR_i\, \varepsilon_{i2} \\
MED_i\, \varepsilon_{i2}
\end{bmatrix}.
\]
Because $\mathbf{x}_{i1}$ and $\mathbf{x}_{i2}$ have the same set of instruments,
each instrument is orthogonal to both $\varepsilon_{i1}$ and $\varepsilon_{i2}$.
\end{quote}

\subsection*{Identification}

Having set up the orthogonality conditions for multiple equations, we can derive
the identification condition in much the same way as in the single-equation case.
The multiple-equation version of (3.3.1) is
\begin{equation}
\mathbf{g}(\mathbf{w}_i;\, \boldsymbol{\delta})
\equiv \begin{bmatrix}
\mathbf{x}_{i1} \cdot (y_{i1} - \mathbf{z}_{i1}' \boldsymbol{\delta}_1) \\
\vdots \\
\mathbf{x}_{iM} \cdot (y_{iM} - \mathbf{z}_{iM}' \boldsymbol{\delta}_M)
\end{bmatrix},
\label{eq:4.1.5}
\end{equation}
where $\mathbf{w}_i$ is as in Assumption~4.2, and $\boldsymbol{\delta}$ without the subscript
is a stacked vector of coefficients:
\begin{equation}
\underset{(\sum_{m=1}^{M} L_m \times 1)}{\boldsymbol{\delta}}
= \begin{bmatrix} \boldsymbol{\delta}_1 \\ \vdots \\ \boldsymbol{\delta}_M \end{bmatrix}.
\label{eq:4.1.6}
\end{equation}

So the orthogonality conditions can be written as
$\E[\mathbf{g}(\mathbf{w}_i;\, \boldsymbol{\delta})] = \mathbf{0}$.
The coefficient vector is identified if $\tilde{\boldsymbol{\delta}} = \boldsymbol{\delta}$
is the only solution to the system of equations
\begin{equation}
\E[\mathbf{g}(\mathbf{w}_i;\, \tilde{\boldsymbol{\delta}})] = \mathbf{0}.
\label{eq:4.1.7}
\end{equation}

Using the definition~\eqref{eq:4.1.5}, we can rewrite the left-hand side of this
equation, $\E[\mathbf{g}(\mathbf{w}_i;\, \tilde{\boldsymbol{\delta}})]$, as
\begin{align}
\E[\mathbf{g}(\mathbf{w}_i;\, \tilde{\boldsymbol{\delta}})]
&\equiv \begin{bmatrix}
\E[\mathbf{x}_{i1} \cdot (y_{i1} - \mathbf{z}_{i1}' \tilde{\boldsymbol{\delta}}_1)] \\
\vdots \\
\E[\mathbf{x}_{iM} \cdot (y_{iM} - \mathbf{z}_{iM}' \tilde{\boldsymbol{\delta}}_M)]
\end{bmatrix} \notag \\
&= \begin{bmatrix}
\E(\mathbf{x}_{i1} \cdot y_{i1}) \\ \vdots \\ \E(\mathbf{x}_{iM} \cdot y_{iM})
\end{bmatrix}
- \begin{bmatrix}
\E(\mathbf{x}_{i1} \mathbf{z}_{i1}') \tilde{\boldsymbol{\delta}}_1 \\
\vdots \\
\E(\mathbf{x}_{iM} \mathbf{z}_{iM}') \tilde{\boldsymbol{\delta}}_M
\end{bmatrix} \notag \\
&= \begin{bmatrix}
\E(\mathbf{x}_{i1} \cdot y_{i1}) \\ \vdots \\ \E(\mathbf{x}_{iM} \cdot y_{iM})
\end{bmatrix}
- \begin{bmatrix}
\E(\mathbf{x}_{i1} \mathbf{z}_{i1}') & & \\
& \ddots & \\
& & \E(\mathbf{x}_{iM} \mathbf{z}_{iM}')
\end{bmatrix}
\begin{bmatrix} \tilde{\boldsymbol{\delta}}_1 \\ \vdots \\ \tilde{\boldsymbol{\delta}}_M \end{bmatrix}
\notag \\
&\equiv \boldsymbol{\sigma}_{\mathbf{x}\mathbf{y}}
- \boldsymbol{\Sigma}_{\mathbf{x}\mathbf{z}}\, \tilde{\boldsymbol{\delta}},
\label{eq:4.1.8}
\end{align}
where
\begin{equation}
\boldsymbol{\sigma}_{\mathbf{x}\mathbf{y}}
\equiv \begin{bmatrix}
\E(\mathbf{x}_{i1} \cdot y_{i1}) \\ \vdots \\ \E(\mathbf{x}_{iM} \cdot y_{iM})
\end{bmatrix}, \quad
\boldsymbol{\Sigma}_{\mathbf{x}\mathbf{z}}
\equiv \begin{bmatrix}
\E(\mathbf{x}_{i1} \mathbf{z}_{i1}') & & \\
& \ddots & \\
& & \E(\mathbf{x}_{iM} \mathbf{z}_{iM}')
\end{bmatrix}.
\label{eq:4.1.9}
\end{equation}
(The third equality in \eqref{eq:4.1.8} is obtained by formula (A.4) of Appendix~A
for multiplication of partitioned matrices.) Therefore, the system of equations
determining $\tilde{\boldsymbol{\delta}}$ can be written as
\begin{equation}
\boldsymbol{\Sigma}_{\mathbf{x}\mathbf{z}}\, \tilde{\boldsymbol{\delta}}
= \boldsymbol{\sigma}_{\mathbf{x}\mathbf{y}},
\label{eq:4.1.10}
\end{equation}
which is the \textit{same} in form as (3.3.4), the system of equations determining
$\tilde{\boldsymbol{\delta}}$ in the single-equation case!

It then follows from the discussion of identification for the single-equation case
that a necessary and sufficient condition for identification is that
$\boldsymbol{\Sigma}_{\mathbf{x}\mathbf{z}}$ be of full column rank.
But, since $\boldsymbol{\Sigma}_{\mathbf{x}\mathbf{z}}$ is block diagonal,
this condition is equivalent to

\begin{assumption}[rank condition for identification]
\label{ass:4.4}
\textit{For each $m$ $(= 1, 2, \ldots, M)$,
$\E(\mathbf{x}_{im} \mathbf{z}_{im}')$ $(K_m \times L_m)$
is of full column rank.}
\end{assumption}

It should be clear why the identification condition becomes thus. We can uniquely
determine all coefficient vectors $(\boldsymbol{\delta}_1, \ldots, \boldsymbol{\delta}_M)$
if and only if each coefficient vector $\boldsymbol{\delta}_m$ is uniquely determined, which occurs
if and only if Assumption~3.4 holds for each equation. The rank condition is this
simple because there are no \textit{a priori} cross-equation restrictions. Identification
when the coefficients are assumed to be the same across equations will be covered
in Section~4.6.

\subsection*{The Assumption for Asymptotic Normality}

As in the single-equation estimation, the orthogonality conditions need to be
strengthened for asymptotic normality:

\begin{assumption}[$\mathbf{g}_i$ is a martingale difference sequence with finite second moments]
\label{ass:4.5}
\textit{$\{\mathbf{g}_i\}$ is a joint martingale difference sequence.
$\E(\mathbf{g}_i \mathbf{g}_i')$ is nonsingular.}
\end{assumption}

The same comment that we made above about joint ergodic stationarity applies
here: the assumption is stronger than requiring the same for each equation. As in
the previous chapters, we use the symbol $\mathbf{S}$ for $\avar(\bar{\mathbf{g}})$,
the asymptotic variance of $\bar{\mathbf{g}}$ (i.e., the variance of the limiting distribution
of $\sqrt{n}\, \bar{\mathbf{g}}$). By the CLT for stationary and ergodic martingale difference
sequences (see Section~2.2), it equals $\E(\mathbf{g}_i \mathbf{g}_i')$.
It has the following partitioned structure:
\begin{equation}
\underset{(\sum_{m=1}^{M} K_m \times \sum_{m=1}^{M} K_m)}{\mathbf{S}}
= \E(\mathbf{g}_i \mathbf{g}_i')
= \begin{bmatrix}
\E(\varepsilon_{i1}\, \varepsilon_{i1}\, \mathbf{x}_{i1} \mathbf{x}_{i1}') & \cdots &
\E(\varepsilon_{i1}\, \varepsilon_{iM}\, \mathbf{x}_{i1} \mathbf{x}_{iM}') \\
\vdots & & \vdots \\
\E(\varepsilon_{iM}\, \varepsilon_{i1}\, \mathbf{x}_{iM} \mathbf{x}_{i1}') & \cdots &
\E(\varepsilon_{iM}\, \varepsilon_{iM}\, \mathbf{x}_{iM} \mathbf{x}_{iM}')
\end{bmatrix}.
\label{eq:4.1.11}
\end{equation}
That is, the $(m, h)$ block of $\mathbf{S}$ is
$\E(\varepsilon_{im}\, \varepsilon_{ih}\, \mathbf{x}_{im} \mathbf{x}_{ih}')$
$(m, h = 1, 2, \ldots, M)$.

In sum, the multiple-equation model is a system of equations where the assumptions
we made for the single-equation model apply to each equation, with the added
requirement of jointness (such as joint stationarity) where applicable.

\subsection*{Connection to the ``Complete'' System of Simultaneous Equations}

The multiple-equation model presented in most other textbooks, called the ``complete''
system of simultaneous equations, adds more assumptions to our model.
Those additional assumptions are \textit{not} needed for the development of multiple-equation
GMM, but if you are curious about them, you can at this junction make a
detour to the first two subsections of Section~8.5.
